\chapter{Perelman's \texorpdfstring{$\mathcal W$}{W} Entropy Functional}

\section{Perelman's \texorpdfstring{$\mathcal W$}{W} entropy functional}

In this chapter, we describe an entropy functional which is more useful than
the $\mathcal{F}$ functional described in Section~6.2. Our main application of
this will be to prove the missing injectivity radius estimate discussed in the
previous chapter.

\subsection{Definition, motivation and basic properties}

Consider now the functional
\[
\mathcal W(g,f,\tau) := \int \bigl[\tau(R + |\nabla f|^2) + f - n\bigr]\,u\,dV,
\]
where $g$ is a metric, $f : \mathcal M \to \mathbb R$ is smooth, $\tau > 0$
is a scale parameter, and $u$ is defined by
\[
u := (4\pi\tau)^{- \frac{n}{2}} e^{-f}.
\]

\begin{definition}
\label{topping-ch8-compatible-arguments}
\dcref{8.1.1}
\group{perelman-w-entropy}
\source{topping:page-0086}
A triple $(g,f,\tau)$ is called \emph{compatible} if
\[
\int_{\mathcal M} u\,dV \equiv \int_{\mathcal M} \frac{e^{-f}}{(4\pi\tau)^{\frac{n}{2}}}\,dV = 1.
\]
\end{definition}

The scale factor $\tau$ is for convenience, in that it could be absorbed into
the other two arguments of $\mathcal W$.

\begin{proposition}
\label{topping-ch8-scale-invariance}
\dcref{8.1.2}
\group{perelman-w-entropy}
\source{topping:page-0087}
\uses{topping-ch8-compatible-arguments}
Under the transformation $(g,f,\tau)\mapsto(\tau^{-1}g,f,1)$, compatibility is preserved, as is the functional $\mathcal W$:
\[
\mathcal W(g,f,\tau)=\mathcal W(\tau^{-1}g,f,1).
\]
\end{proposition}

One motivation for the $\mathcal W$-entropy comes from the log-Sobolev
inequality of Gross, as we now describe. We will later show how $\mathcal W$
arises from considering the classical entropy from Section~6.5.

On $\mathbb R^n$, the Gaussian measure $d\mu$ is defined in terms of the
Lebesgue measure $dx$ by
\[
d\mu = (2\pi)^{-\frac n2} e^{-\frac{|x|^2}{2}}\,dx,
\]
the normalisation being chosen so that
\[
\int_{\mathbb R^n} d\mu = 1.
\]

\begin{theorem}
\label{topping-ch8-gross-log-sobolev}
\dcref{8.1.3}
\group{perelman-w-entropy}
\source{topping:page-0087}
If $v:\mathbb R^n\to\mathbb R$ is, say, smooth and satisfies $v, |\nabla v|\in L^2(d\mu)$ then
\[
\int v^2\ln|v|\,d\mu \le \int |\nabla v|^2\,d\mu + \left(\int v^2\,d\mu\right) \ln\left(\int v^2\,d\mu\right)^{\frac12}.
\]
\end{theorem}

\begin{remark}
\label{topping-ch8-gross-equality}
\dcref{8.1.4}
\group{perelman-w-entropy}
\source{topping:page-0087}
We have equality in the preceding inequality if and only if $v$ is constant.
\end{remark}

\begin{remark}
\label{topping-ch8-gross-no-dimension-constant}
\dcref{8.1.5}
\group{perelman-w-entropy}
\source{topping:page-0087}
There is no constant in this inequality which depends on the dimension $n$.
\end{remark}

\begin{remark}
\label{topping-ch8-gross-unit-mass}
\dcref{8.1.6}
\group{perelman-w-entropy}
\source{topping:page-0087}
If we choose $v$ so that $\int v^2\,d\mu=1$, then the inequality becomes
\[
\int v^2\ln|v|\,d\mu \le \int |\nabla v|^2\,d\mu.
\]
\end{remark}

Let us consider the consequences of this inequality for the functional
$\mathcal W$ in the case that $(\mathcal M,g)$ is Euclidean space.

\begin{lemma}
\label{topping-ch8-euclidean-w-nonnegative}
\dcref{8.1.7}
\group{perelman-w-entropy}
\source{topping:page-0087,topping:page-0088,topping:page-0089}
\uses{topping-ch8-gross-log-sobolev,topping-ch8-scale-invariance}
Let $g$ denote the flat metric on $\mathcal M=\mathbb R^n$. If $f$ and $\tau$
are compatible with $g$, then
\[
\mathcal W(g,f,\tau)\ge 0
\]
with equality if and only if $f(x)\equiv \frac{|x|^2}{4\tau}$.
\end{lemma}

\begin{proof}
For now, set $\tau=\frac12$ and let $f:\mathbb R^n\to\mathbb R$ be compatible
with $g$ and $\tau$, which in this situation means that
\[
\int_{\mathbb R^n} \frac{e^{-f}}{(2\pi)^{\frac n2}}\,dx = 1.
\]

If we define $v=e^{\frac{|x|^2}{4}-\frac f2}$, we have
\[
v^2\,d\mu
= e^{\frac{|x|^2}{2}-f}(2\pi)^{-\frac n2}e^{-\frac{|x|^2}{2}}\,dx
= (2\pi)^{-\frac n2}e^{-f}\,dx
\]
and so $\int_{\mathbb R^n} v^2\,d\mu=1$. Therefore, by the log-Sobolev
inequality,
\[
\int v^2\ln|v|\,d\mu \le \int |\nabla v|^2\,d\mu.
\]

We wish to convert from $v$ to $f$ in this inequality. For the left-hand side,
we have
\[
\int v^2\ln|v|\,d\mu
= \int e^{\frac{|x|^2}{2}-f}\left(\frac{|x|^2}{4}-\frac f2\right)
(2\pi)^{-\frac n2}e^{-\frac{|x|^2}{2}}\,dx
\]
\[
= \int \left(\frac{|x|^2}{4}-\frac f2\right)\frac{e^{-f}}{(2\pi)^{\frac n2}}\,dx.
\]

For the right-hand side, first note that $\nabla\cdot x=n$ to give the
integration-by-parts formula
\[
-\int \frac{x\cdot\nabla f}{2}\,\frac{e^{-f}}{(2\pi)^{\frac n2}}\,dx
= \frac12 \int x\cdot\nabla\!\left(e^{-f}\right)\frac{dx}{(2\pi)^{\frac n2}}
= -\frac n2 \int \frac{e^{-f}}{(2\pi)^{\frac n2}}\,dx.
\]

Therefore,
\[
\int |\nabla v|^2\,d\mu
= \int \left(\frac{|x|^2}{4}-\frac n2+\frac{|\nabla f|^2}{4}\right)
\frac{e^{-f}}{(2\pi)^{\frac n2}}\,dx,
\]
and the log-Sobolev inequality becomes
\[
\int \left(\frac{|x|^2}{4} - \frac{f}{2}\right)\frac{e^{-f}}{(2\pi)^{\frac{n}{2}}}\,dx
\leq
\int \left(\frac{|x|^2}{4} - \frac{n}{2} + \frac{|\nabla f|^2}{4}\right)\frac{e^{-f}}{(2\pi)^{\frac{n}{2}}}\,dx.
\]
Equivalently,
\[
\mathcal W\!\left(g,f,\frac{1}{2}\right)
:=
\int \left[\frac{1}{2}|\nabla f|^2 + f - n\right]\frac{e^{-f}}{(2\pi)^{\frac{n}{2}}}\,dx
\geq 0.
\]
By the scale invariance of $\mathcal W$ and the fact that $\mathbb R^n$ is
preserved under homothetic scaling, we conclude that
\[
\mathcal W(g,f,\tau)\geq 0
\]
for any $f,\tau$ compatible with $g$, with equality if and only if
$f(x)\equiv \frac{|x|^2}{4\tau}$.
\end{proof}

In contrast to the situation for the $\mathcal F$-functional in Section~6.6,
one must take a little care to show that the $\mathcal W$-functional is
bounded below and that a minimising compatible $f$ exists, given $g$ and
$\tau$.

\begin{lemma}
\label{topping-ch8-minimiser-exists}
\dcref{8.1.8}
\group{perelman-w-entropy}
\source{topping:page-0089,topping:page-0090,topping:page-0091}
\uses{topping-ch8-scale-invariance}
For any smooth Riemannian metric $g$ on a closed manifold $\mathcal M$, and
$\tau>0$, the infimum of $\mathcal W(g,f,\tau)$ over all compatible $f$ (as
defined in Definition~8.1.1) is attained by a smooth compatible $f$.
\end{lemma}

Defining $\mu$ as this infimum,
\[
\mu(g,\tau):=\inf_f \mathcal W(g,f,\tau),
\]
the function $\mu(g,\tau)$ is bounded below as $\tau$ varies within any finite
interval $(0,\tau_0]$.

Given this lemma, we may define, for $\tau_0>0$,
\[
\nu(g,\tau_0):=\inf_{\tau\in(0,\tau_0]}\mu(g,\tau)>-\infty.
\]

The starting point for the proof of this lemma is the change of variables
$\phi:=e^{-f/2}$ that we used in Section~6.6. Abusing notation for
$\mathcal W$, we then have
\[
\mathcal W(g,\phi,\tau)
=(4\pi\tau)^{-\frac{n}{2}}
\int \left[\tau\!\left(4|\nabla\phi|^2+\Ric\,\phi^2\right)-2\phi^2\ln\phi-n\phi^2\right]\,dV,
\]
and the compatibility constraint becomes
\[
(4\pi\tau)^{-\frac{n}{2}} \int \phi^2\,dV = 1.
\]

These expressions have the benefit of making perfect sense when $\phi$ is
merely weakly (rather than strictly) positive, and by approximation,
\[
\inf_f \mathcal W(g,f,\tau) = \inf_\phi \mathcal W(g,\phi,\tau),
\]
where the infima are taken over $f : \mathcal M \to \mathbb R$ and
$\phi : \mathcal M \to [0,\infty)$ compatible with $g$ and $\tau$, and we are
abusing notation for $\mathcal W$ as usual.

We will use this change of variables to prove the existence of a minimising
compatible $\phi$, given $g$ and $\tau$. One can argue that this minimiser will
be strictly positive and hence one has the existence of a minimising $f$ as
asserted in Lemma~8.1.8.

We will not prove the asserted boundedness of $\mu$, which is from [31], since
it requires a little more analysis than we are assuming.

\begin{proof}
(Parts of Lemma~8.1.8.) We limit ourselves to proving that the functional
$\mathcal W$ is bounded below, for given $g$ and $\tau$, and that a minimising
sequence of compatible positive functions $\phi$ must be bounded in $W^{1,2}$,
and hence that the direct method applies to give a minimiser. (See [15, \S8.12].)

By the invariance observed in Proposition~8.1.2, we may scale $g$ and $\tau$
simultaneously and assume that $\tau = \frac{1}{4\pi}$, say, which slightly
simplifies the preceding formulas.

It suffices to prove that
\[
\mathcal W(\phi) := \int \left[ \frac{1}{4\pi}\left(4|\nabla \phi|^2 + \Ric\, \phi^2\right) - 2\phi^2 \ln \phi - n\phi^2 \right] dV
\]
\[
\geq \frac{1}{2\pi} \int |\nabla \phi|^2 dV - C(g,n),
\]
for any positive $\phi$ satisfying
\[
\int \phi^2\,dV = 1,
\]
for some constant $C(g,n)$.

To achieve this, we split up the expression for $\mathcal W(\phi)$. First, by
the normalisation we clearly have
\[
\int \left[\frac{1}{4\pi}(R\phi^2)-n\phi^2\right]\,dV \ge -C(g,n),
\]
since our manifold is closed, so $R$ is bounded.

Meanwhile, we may write
\[
\int \phi^2 \ln \phi\,dV
= \frac{n-2}{4}\int \phi^2 \ln \phi^{\frac{4}{n-2}}\,dV
= \frac{n-2}{4}\int G(\sigma)\,d\mu
\]
where $\sigma:\mathcal M\to (0,\infty)$ is defined by $\sigma := \phi^{\frac{4}{n-2}}$, the function $G:(0,\infty)\to \mathbb R$ defined by $G(y):=\ln y$ is concave, and the measure $d\mu := \phi^2\,dV$ has unit total mass. We may apply Jensen's inequality to tell us that
\[
\int G(\sigma)\,d\mu \le G\left(\int \sigma\,d\mu\right),
\]
and therefore
\[
\int \phi^2 \ln \phi\,dV
\le \frac{n-2}{4}\ln\left[\int \phi^{2+\frac{4}{n-2}}\,dV\right].
\]
We continue by applying the Sobolev inequality to give
\[
\int \phi^2 \ln \phi\,dV \le \frac{1}{2\pi}\int |\nabla\phi|^2\,dV + C(g,n),
\]
where the constant $C(g,n)$ is permitted to change at each step.

Combining with the earlier estimate, we conclude that
\[
\mathcal W(\phi) \ge \frac{1}{2\pi}\int |\nabla\phi|^2\,dV - C(g,n),
\]
as required.
\end{proof}

\subsection{Monotonicity of $\mathcal W$}

In this section, we would like to show that $\mathcal W$, like
$\mathcal F$, is increasing under the Ricci flow when $f$ and $\tau$ are made
to evolve appropriately.

\begin{proposition}
\label{topping-ch8-monotonicity}
\dcref{8.2.1}
\group{perelman-w-entropy}
\source{topping:page-0092}
If $\mathcal M$ is closed, and $g$, $f$ and $\tau$ evolve according to
\[
\begin{cases}
\displaystyle \frac{\partial g}{\partial t} = -2\Ric \\
\displaystyle \frac{d\tau}{dt} = -1 \\
\displaystyle \frac{\partial f}{\partial t} = -\Delta f + \lvert \nabla f\rvert^{2} - R + \frac{n}{2\tau}
\end{cases}
\]
then the functional $\mathcal W$ increases according to
\[
\frac{d}{dt}\mathcal W(g,f,\tau) = 2\tau \int \lvert \Ric + \Hess(f) - \frac{g}{2\tau}\rvert^{2} \,dV \geq 0.
\]
\end{proposition}

\begin{remark}
\label{topping-ch8-compatibility-preserved}
\dcref{8.2.2}
\group{perelman-w-entropy}
\source{topping:page-0092}
Under the evolution of the previous proposition, $u$ satisfies
\[
\Box^{*}u = 0,
\]
and the compatibility constraint is preserved.
\end{remark}

\begin{remark}
\label{topping-ch8-shrinking-soliton}
\dcref{8.2.3}
\group{perelman-w-entropy}
\source{topping:page-0092}
If the time derivative of $\mathcal W$ in the monotonicity formula ever fails to
be strictly positive, then $\Ric + \Hess(f) - \frac{g}{2\tau} = 0$, and the
flow is a shrinking gradient soliton.
\end{remark}

\begin{remark}
\label{topping-ch8-classical-entropy-link}
\dcref{8.2.4}
\group{perelman-w-entropy}
\source{topping:page-0092}
The $\mathcal W$-functional applied to this backwards Brownian diffusion on a
Ricci flow also arises via the renormalised classical entropy $\tilde N$.
\end{remark}

\begin{remark}
\label{topping-ch8-existence-of-coupled-system}
\dcref{8.2.5}
\group{perelman-w-entropy}
\source{topping:page-0093}
As for the existence of a solution to the coupled system, once we have a Ricci
flow on a time interval $[0,T]$ we may set $\tau = T + C - t$ and solve for
$f$ with prescribed final data $f(T)$ by considering the linear equation
satisfied by $u$.
\end{remark}

The following result can be considered to be a local version of
Proposition~8.2.1.

\begin{proposition}
\label{topping-ch8-local-monotonicity}
\dcref{8.2.6}
\group{perelman-w-entropy}
\source{topping:page-0093,topping:page-0094,topping:page-0095}
\uses{topping-ch8-monotonicity,topping-ch8-compatibility-preserved}
Suppose $g$, $f$ and $\tau$ evolve as in Proposition~8.2.1, and $u$ is
defined as in the definition of $\mathcal W$. Then the function
\[
v := \bigl[\tau(2\Delta f - |\nabla f|^{2} + R) + f - n\bigr]u
\]
satisfies
\[
\Box^{*} v = -2\tau\lvert \Ric + \Hess(f) - \tfrac{g}{2\tau}\rvert^{2}u.
\]
\end{proposition}

\begin{remark}
\label{topping-ch8-proposition-implies-monotonicity}
\dcref{8.2.7}
\group{perelman-w-entropy}
\source{topping:page-0093}
Proposition~8.2.1 follows immediately from Proposition~8.2.6 because
\[
\mathcal W = \int_{\mathcal M} v\,dV.
\]
\end{remark}

\begin{proof}
(Proposition~8.2.6.) Splitting $v$ as a product, we find that
\[
\Box^{*} v = \Box^{*}\!\left(\frac{v}{u}u\right)
= \frac{v}{u}\,\Box^{*} u - u\!\left(\frac{\partial}{\partial t} + \Delta\right)\!\left(\frac{v}{u}\right)
- 2\langle \nabla\!\left(\frac{v}{u}\right), \nabla u\rangle,
\]
and since $\Box^{*}u = 0$, we have
\[
u^{-1}\Box^{*} v = -\left(\frac{\partial}{\partial t} + \Delta\right)\!\left(\frac{v}{u}\right)
- 2u^{-1}\langle \nabla\!\left(\frac{v}{u}\right), \nabla u\rangle.
\]

For the first of the terms on the right-hand side, we compute
\[
-\Bigl(\frac{\partial}{\partial t}+\Delta\Bigr)\Bigl(\frac vu\Bigr)
=-\Bigl(\frac{\partial}{\partial t}+\Delta\Bigr)\bigl[\tau(2\Delta f-|\nabla f|^{2}+R)+f-n\bigr]
\]
\[
=\bigl(2\Delta f-|\nabla f|^{2}+R\bigr)
-\tau\Bigl(\frac{\partial}{\partial t}+\Delta\Bigr)\bigl(2\Delta f-|\nabla f|^{2}+R\bigr)
-\Bigl(\frac{\partial}{\partial t}+\Delta\Bigr)f,
\]
and using the evolution equation for $f$ this reduces to
\[
-\Bigl(\frac{\partial}{\partial t}+\Delta\Bigr)\Bigl(\frac vu\Bigr)
=2\Delta f-2|\nabla f|^{2}+2R-\frac{n}{2\tau}
-\tau\Bigl(\frac{\partial}{\partial t}+\Delta\Bigr)\bigl(2\Delta f-|\nabla f|^{2}+R\bigr).
\]

We deal with each part of the final term individually. By the evolution of
$f$,
\[
\Bigl(\frac{\partial}{\partial t}+\Delta\Bigr)\Delta f
=2\langle \Ric,\Hess(f)\rangle+\Delta\Bigl(|\nabla f|^{2}-R+\frac{n}{2\tau}\Bigr).
\]
Using also the evolution equation for $g$ and the evolution of $R$,
\[
\Bigl(\frac{\partial}{\partial t}+\Delta\Bigr)\bigl(2\Delta f-|\nabla f|^{2}+R\bigr)
=4\langle \Ric,\Hess(f)\rangle+\Delta|\nabla f|^{2}-2\Ric(\nabla f,\nabla f)
\]
\[
-2\langle \nabla f,\nabla(-\Delta f+|\nabla f|^{2}-R)\rangle+2|\Ric|^{2}.
\]

Also,
\[
\Bigl\langle \nabla\frac vu,\nabla u\Bigr\rangle
=\bigl\langle \nabla\bigl(\tau(2\Delta f-|\nabla f|^{2}+R)+f\bigr),-u\nabla f\bigr\rangle,
\]
so
\[
-2u^{-1}\Bigl\langle\nabla\frac vu,\nabla u\Bigr\rangle
=2\tau\Bigl\langle\nabla\bigl(2\Delta f-|\nabla f|^{2}+R\bigr),\nabla f\Bigr\rangle+2|\nabla f|^{2}.
\]

Combining these expressions, we find that
\[
u^{-1}\Box^{*}v
=-\frac{n}{2\tau}+2\langle \Ric+\Hess(f),g\rangle-2\tau|\Ric+\Hess(f)|^{2}
\]
\[
=-2\tau\Bigl|\Ric+\Hess(f)-\frac{g}{2\tau}\Bigr|^{2}.
\]
\end{proof}

\subsection{No local volume collapse where curvature is controlled}

Given a complete Riemannian manifold $(\mathcal M,g)$, a point $p\in\mathcal M$
and $s>0$, we use the shorthand
\[
\mathcal V(p,s)=\mathcal V_g(p,s):=\Vol\bigl(B_g(p,s)\bigr)
\]
for the volume of the geodesic ball $B(p,s)=B_g(p,s)$ centred at $p$ of
radius $s$, computed with respect to the metric $g$. Our aim is to get a lower
bound on the volume ratio
\[
\mathcal K(p,r):=\frac{\mathcal V(p,r)}{r^n},
\]
during a Ricci flow which will be turned into a lower bound for the injectivity
radius, in terms of the maximum of the curvature $|{\rm Rm}|$, in the next
section.

\begin{theorem}
\label{topping-ch8-no-collapsing}
\dcref{8.3.1}
\group{perelman-w-entropy}
\source{topping:page-0096}
\uses{topping-ch8-mu-ball-estimate}
Suppose that $g(t)$ is a Ricci flow on a closed manifold $\mathcal M$, for
$t \in [0,T]$. Working with respect to the metric $g(T)$, if $p \in \mathcal M$,
and $r > 0$ is sufficiently small so that $|R| \le r^{-2}$ on $B(p,r)$, then
\[
\mathcal K(p,r) > \xi
\]
for some $\xi > 0$ depending on $n$, $g(0)$, and upper bounds for $r$ and $T$.
\end{theorem}

\begin{remark}
\label{topping-ch8-no-collapsing-variant}
\dcref{8.3.2}
\group{perelman-w-entropy}
\source{topping:page-0096}
We could replace the $|R| \le r^{-2}$ hypothesis by $|R| \le C r^{-2}$,
although $\xi$ would then also depend on $C$.
\end{remark}

\begin{remark}
\label{topping-ch8-no-collapsing-perelman}
\dcref{8.3.3}
\group{perelman-w-entropy}
\source{topping:page-0096}
This is modelled on Perelman's no collapsing theorem from [31, \S4].
\end{remark}

We will prove the theorem above by combining an iteration argument with the
theorem below.

\begin{theorem}
\label{topping-ch8-mu-ball-estimate}
\dcref{8.3.4}
\group{perelman-w-entropy}
\source{topping:page-0096}
\uses{topping-ch8-minimiser-exists,topping-ch8-monotonicity}
Suppose that $g(t)$ is a Ricci flow on a closed manifold $\mathcal M$, for
$t \in [0,T]$, and that $r > 0$ and $p \in \mathcal M$. Then computing
volumes, curvature and geodesic balls with respect to $g(T)$, we have
\[
\gamma \le \frac{\mathcal V(p,r)}{\mathcal V(p,\frac r2)}
+ \frac{r^2}{\mathcal V(p,\frac r2)} \int_{B(p,r)} |R|\,dV + \ln[\mathcal K(p,r)],
\]
for some $\gamma \in \mathbb R$ depending on $n$, $g(0)$, and upper bounds for
$r$ and $T$.
\end{theorem}

In the proof of this theorem, the monotonicity of $\mathcal W$ will give a
lower bound for $\mathcal W$ under the flow, which will then be turned into
geometric information via the following lemma.

\begin{lemma}
\label{topping-ch8-mu-ball-bound}
\dcref{8.3.5}
\group{perelman-w-entropy}
\source{topping:page-0096,topping:page-0097,topping:page-0098}
\uses{topping-ch8-minimiser-exists}
For any smooth metric $g$ on a closed manifold $\mathcal M$, and $r > 0$,
$p \in \mathcal M$ and $\lambda > 0$,
\[
\mu(g,\lambda r^2) \le 36\lambda \frac{\mathcal V(p,r) - \mathcal V(p,\frac r2)}{\mathcal V(p,\frac r2)}
+ \frac{\lambda r^2}{\mathcal V(p,\frac r2)} \int_{B(p,r)} |R|\,dV + \ln\!\Bigl[\frac{\mathcal V(p,r)}{(4\pi \lambda r^2)^{n/2}}\Bigr] - n.
\]
\end{lemma}

\begin{proof}
(Lemma~8.3.5.) If we adopt the change of variables $\phi = e^{-f/2}$ and write
$\tau = \lambda r^2$, then we have that
\[
\mathcal W(g,\phi,\lambda r^2)
= (4\pi \lambda r^2)^{-n/2}
\int \Bigl[\lambda r^2(4|\nabla \phi|^2 + R\phi^2) - 2\phi^2 \ln \phi - n\phi^2\Bigr]\,dV,
\]
and the compatibility constraint becomes
\[
(4\pi \lambda r^2)^{-n/2}\int \phi^2\,dV = 1.
\]

As we remarked before, these expressions have the benefit of making sense when
$\phi$ is merely weakly positive, and by approximation,
\[
\inf_f \mathcal W(g,f,\lambda r^2)
= \inf_\phi \mathcal W(g,\phi,\lambda r^2),
\]
where the infima are taken over compatible $f$ and $\phi$.

Let $\psi : [0,\infty) \to [0,1]$ be a smooth cut-off function, supported in
$[0,1]$, such that $\psi(y)=1$ for $y \in [0,\tfrac12]$ and $|\psi'| \le 3$.
We then write
\[
\phi(x) = e^{-c/2}\psi\!\left(\frac{d(x,p)}{r}\right),
\]
where $c \in \mathbb R$ is determined by the constraint, and since
\[
\mathcal V(p,\tfrac r2) \le e^c \int \phi^2\,dV \le \mathcal V(p,r),
\]
we deduce that
\[
(4\pi \lambda r^2)^{-n/2}\mathcal V(p,\tfrac r2)
\le e^c \le
(4\pi \lambda r^2)^{-n/2}\mathcal V(p,r).
\]

We now estimate each of the four terms separately.

For the gradient term, whose gradient is supported on $B(p,r)\setminus
B(p,\tfrac r2)$ and satisfies $|\nabla \phi|\le \frac{3}{r}e^{-c/2}$, we
estimate
\[
(4\pi \lambda r^2)^{-n/2}
\int \lambda r^2 4|\nabla \phi|^2\,dV
\le 36\lambda \frac{\mathcal V(p,r)-\mathcal V(p,\tfrac r2)}{\mathcal V(p,\tfrac r2)}.
\]

For the scalar curvature term, since $\phi$ is supported on $B(p,r)$ and
$\phi^2 \le e^{-c} \le \frac{(4\pi \lambda r^2)^{\frac{n}{2}}}{\mathcal V(p,\frac{r}{2})}$, we estimate
\[
(4\pi \lambda r^2)^{-\frac{n}{2}} \int \lambda r^2 R \phi^2\,dV \le \frac{\lambda r^2}{\mathcal V(p,\frac{r}{2})} \int_{B(p,r)} |R|\,dV.
\]

For the logarithmic term, we rewrite
\[
(4\pi \lambda r^2)^{-\frac{n}{2}} \int -2\phi^2 \ln \phi\,dV = \int G(\sigma)\,d\mu,
\]
where $\sigma := \phi^2$, $G(y) := -y\ln y$, and $d\mu := (4\pi \lambda r^2)^{-\frac n2}dV$.
Because $G$ is concave, Jensen's inequality gives
\[
\int G(\sigma)\,d\mu \le G\!\left(\int \sigma\,d\mu\right),
\]
and since the normalisation gives $\int \sigma\,d\mu = 1$, this implies
\[
(4\pi \lambda r^2)^{-\frac{n}{2}} \int -2\phi^2 \ln \phi\,dV \le \ln\!\left[\frac{\mathcal V(p,r)}{(4\pi \lambda r^2)^{\frac{n}{2}}}\right].
\]

For the final term, the constraint gives simply
\[
(4\pi \lambda r^2)^{-\frac{n}{2}} \int -n\phi^2\,dV = -n.
\]

Combining these calculations gives the claimed estimate.
\end{proof}

\begin{proof}
(Theorem~8.3.4.) First, let us specialise Lemma~8.3.5 to the case
\[
\lambda = \frac{1}{36}\quad\text{and}\quad g=g(T),
\]
and estimate
\[
\mu\!\left(g(T),\frac{1}{36}r^2\right)
\le \frac{\mathcal V(p,r)}{\mathcal V(p,\frac{r}{2})}
+\frac{r^2}{\mathcal V(p,\frac{r}{2})}\int_{B(p,r)} |R|\,dV
+\ln[\mathcal K(p,r)]-\frac{n}{2}\ln\frac{\pi}{9}-n.
\]

By Lemma~8.1.8, there exists a smooth $f_T : \mathcal M \to \mathbb R$
compatible with $g(T)$ and $\tau=\frac{1}{36}r^2$ such that
\[
\mathcal W\!\left(g(T),f_T,\frac{1}{36}r^2\right)
=\mu\!\left(g(T),\frac{1}{36}r^2\right).
\]

We set $\tau = T + \frac{1}{36}r^2 - t$, and for our given Ricci flow $g(t)$,
find the $f : \mathcal M\times[0,T] \to \mathbb R$ with $f(T)=f_T$
completing a solution of the coupled system, as discussed earlier. By the
compatibility preservation remark, $g(t)$, $f(t)$ and $\tau$ remain compatible
for all $t$. Using the definition of $\mu$, Proposition~8.2.1 and the
preceding identity we then have
\[
\mu\!\left(g(0),\frac{1}{36}r^2+T\right)
\le \mathcal W\!\left(g(0),f(0),\frac{1}{36}r^2+T\right)
 \le \mathcal W\!\left(g(T),f(T),\frac{1}{36}r^2\right)
= \mu\!\left(g(T),\frac{1}{36}r^2\right).
\]
This, coupled with the earlier estimate and the definition of $\nu$, gives the
claimed inequality.
\end{proof}

\begin{proof}
(Theorem~8.3.1.) Let $r_0$ and $T_0$ be upper bounds for $r$ and $T$
respectively. By Theorem~8.3.4, there exists $\gamma=\gamma(n,g(0),T_0,r_0)\in\mathbb R$ such that
working with respect to the metric $g(T)$, we have, for all $s\in(0,r_0]$,
that
\[
\gamma \le \frac{\mathcal V(p,s)}{\mathcal V(p,\frac{s}{2})}
+\frac{s^2}{\mathcal V(p,\frac{s}{2})}\int_{B(p,s)} |R|\,dV + \ln[\mathcal K(p,s)].
\]

When $s \in (0,r]$, we have that $|R| \le r^{-2} \le s^{-2}$ on $B(p,s)$, and
therefore
\[
\gamma \le 2 \frac{\mathcal V(p,s)}{\mathcal V(p,\frac{s}{2})} + \ln [\mathcal K(p,s)].
\]

Define $\omega_n$ to be the volume of the unit ball in Euclidean $n$-space.
Since $g(T)$ is a smooth metric, we have
\[
\mathcal K(p,s) = \frac{\mathcal V(p,s)}{s^n} \to \omega_n
\]
as $s \downarrow 0$. We claim that the theorem holds true if we choose
\[
\xi := \min\left\{ \frac{\omega_n}{2}, e^{\gamma - 2^{n+1}} \right\} > 0.
\]
That is, we wish to establish that $\mathcal K(p,r) > \xi$.

Claim: If $s \in (0,r]$ and $\mathcal K(p,s) \le \xi$ then $\mathcal K(p,\frac{s}{2}) \le \xi$.

Indeed if $\mathcal K(p,s) \le \xi$, then $\mathcal K(p,s) \le e^{\gamma - 2^{n+1}}$, so
$\ln[\mathcal K(p,s)] \le \gamma - 2^{n+1}$, and by the previous inequality,
\[
\gamma \le 2 \frac{\mathcal V(p,s)}{\mathcal V(p,\frac{s}{2})} + \gamma - 2^{n+1}.
\]
Rearranging, we find that
\[
2^n \le \frac{\mathcal V(p,s)}{\mathcal V(p,\frac{s}{2})} = 2^n \frac{\mathcal K(p,s)}{\mathcal K(p,\frac{s}{2})},
\]
and hence that $\mathcal K(p,\frac{s}{2}) \le \mathcal K(p,s) \le \xi$, as
required for the claim.

Using this claim iteratively, we see that
\[
\mathcal K(p,r) \le \xi \qquad \Longrightarrow \qquad \mathcal K(p,2^{-m}r) \le \xi \le \frac{\omega_n}{2}
\]
for all $m \ge 1$, but this contradicts the limit
$\lim_{s\downarrow 0}\mathcal K(p,s)=\omega_n$. Therefore $\mathcal K(p,r)>\xi$.
\end{proof}

\subsection{Volume ratio bounds imply injectivity radius bounds}

The output of Theorem~8.3.1 is a positive lower bound on the volume ratio
$\mathcal K$ at scales no larger than the inverse of the square-root of the
curvature. In this section, we show that this implies a positive lower bound
on the injectivity radius, as required to complete the discussion of
singularity blow-ups that we began in Section~7.3.

Given a Riemannian manifold $(\mathcal M,g)$, we denote by
\[
\operatorname{inj}(\mathcal M) = \operatorname{inj}(\mathcal M, g) := \inf_{p \in \mathcal M} \operatorname{inj}(\mathcal M, g, p),
\]
its injectivity radius.

\begin{lemma}
\label{topping-ch8-volume-ratio-inj-estimate}
\dcref{8.4.1}
\group{perelman-w-entropy}
\source{topping:page-0101,topping:page-0102,topping:page-0103}
\uses{topping-ch8-no-collapsing}
There exist $\bar r > 0$ universal and $\eta > 0$ depending on the dimension
$n$ such that if $(\mathcal M^n, g)$ is a closed Riemannian manifold
satisfying $|\operatorname{Rm}| \le 1$, then there exists $p \in \mathcal M$
such that
\[
\frac{\mathcal V(p,r)}{r^n} \le \frac{\eta}{r}\,\operatorname{inj}(\mathcal M),
\]
for all $r \in (0,\bar r]$.
\end{lemma}

In order to see the sharpness of this lemma, one can consider $\mathcal M$ to
be the product of a closed manifold of dimension $n-1$ satisfying
$|\operatorname{Rm}| \le 1$, and a very small $S^1$.

\begin{proof}
For the moment, set $\bar r = 1$. We will arrive at the constant $\bar r$
whose existence is asserted in the proof by successively reducing it to satisfy
a number of constraints.

First we observe that it suffices to prove the theorem under the assumption
that $\operatorname{inj}(\mathcal M)$ is less than $\pi$. This follows because
our curvature hypothesis $|\operatorname{Rm}| \le 1$ implies that all
sectional curvatures are no lower than $-1$, and hence by Bishop's theorem the
volume ratio on the left-hand side is bounded above for $r \le 1$ by some
$n$-dependent constant $C$. Whenever $\operatorname{inj}(\mathcal M)\ge \pi$,
the theorem holds with $\eta = \frac{C}{\pi}$.

We now appeal to Klingenberg's lemma, which tells us that because all of the
sectional curvatures are no more than 1, and $\operatorname{inj}(\mathcal M)<\pi$,
there exists a closed geodesic loop $\gamma$ in $\mathcal M$ with
\[
\operatorname{Length}(\gamma) = 2\,\operatorname{inj}(\mathcal M).
\]

We are free to pick any $p \in \gamma$.

Let us define $N^r\gamma$ to be the subset of the normal bundle of $\gamma$
within $\mathcal M$ consisting of vectors of length less than $r$, and
$\mathcal N^r\gamma$ to be the subset of $\mathcal M$ consisting of all points
within a distance $r$ of $\gamma$. There is a natural flat metric $G$ on the
total space of $N^r\gamma$, defined unambiguously at points over a sufficiently
small simply connected neighbourhood $U\subset \gamma$ as the pullback of the
standard product metric on $\gamma \times N_q^r\gamma$.

We define a map $u : N^r\gamma \to \mathcal N^r\gamma$ by exponentiation. That
is, for $v \in N^r\gamma$,
\[
u(v) = \exp_{\pi(v)}(v).
\]

Clearly, $u$ is a surjection. Moreover, by considering the behaviour of Jacobi
fields under our curvature bound and reducing $\bar r$ if necessary, the map
$u$ is an immersion and we have
\[
u^*g \le 4G,
\]
say. In particular, the volume form of $u^*g$ is bounded by $2^n$ times the
volume form of $G$, and hence
\[
\Vol_g(\mathcal N^r\gamma) \le 2^n \operatorname{Length}(\gamma)\omega_{n-1}r^{n-1},
\]
where $\omega_{n-1}$ is the volume of the $(n-1)$-dimensional Euclidean unit
ball.

Clearly $B_g(p,r) \subset \mathcal N^r\gamma$, and so
\[
\mathcal V(p,r) \le \Vol(\mathcal N^r\gamma) \le 2^n \omega_{n-1}\operatorname{Length}(\gamma)r^{n-1},
\]
and by the geodesic-loop identity we have
\[
\frac{\mathcal V(p,r)}{r^n} \le 2^{n+1}\,\omega_{n-1}\frac{\operatorname{inj}(\mathcal M)}{r}.
\]
\end{proof}

\subsection{Blowing up at singularities II}

Now we have the no collapsing result of Theorem~8.3.1, and the injectivity
radius control of Lemma~8.4.1, we can continue the discussion of blowing up
Ricci flows near singularities that we began in Section~7.3.

Recall that for a Ricci flow $g(t)$ on a closed manifold $\mathcal M$, over a
maximal time interval $[0,T)$ with $T<\infty$, we found appropriate sequences
$\{p_i\}\subset \mathcal M$ and $t_i \uparrow T$, with
\[
|\Rm|(p_i,t_i)=\sup_{x\in \mathcal M,\, t\in [0,t_i]} |\Rm|(x,t) \to \infty,
\]
and defined blown-up Ricci flows
\[
g_i(t):=|\Rm|(p_i,t_i)\, g\!\left(t_i+\frac{t}{|\Rm|(p_i,t_i)}\right)
\]
such that, for all $a<0$ and some $b=b(n)>0$, the curvature of $g_i(t)$ was
controlled for $t\in (a,b)$, and in particular
\[
\sup_{x\in \mathcal M,\, t\in (a,b)} |R(g_i(t))| \le M < \infty,
\]
for sufficiently large $i$, and $M$ dependent only on the dimension $n$. By
choosing $r\in (0,M^{-1/2})$, we have $|R(g_i(0))|\le r^{-2}$. Converting back
to the original flow, before blowing up, we see that if
$0<r\le (M|\Rm|(p_i,t_i))^{-1/2}$, then $|R(g(t_i))|\le r^{-2}$.

By Theorem~8.3.1 we know that for all $p\in \mathcal M$,
$0<r\le (M|\Rm|(p_i,t_i))^{-1/2}$ and sufficiently large $i$, we have the
lower bound
\[
\frac{\mathcal V(p,r)}{r^n}>\xi>0
\]
at time $t=t_i$, where $\xi=\xi(n,g(0),T)>0$. Here, we included the condition
``for sufficiently large $i$'' so that we are free to ignore the dependency of
$\xi$ on an upper bound for $r$ in Theorem~8.3.1.

Returning to the blown-up flows $g_i(t)$, we find that for all $p \in \mathcal M$,
with respect to the metric $g_i(0)$,
\[
\frac{\mathcal V(p,r)}{r^n} > \xi,
\]
for all $r \in (0, M^{-1/2}]$.

By Lemma~8.4.1, after fixing $r = \min\{M^{-1/2}, \bar r\} > 0$ this implies
the lower bound on the injectivity radius,
\[
\operatorname{inj}(\mathcal M,g_i(0)) \ge \frac{r}{\eta}\frac{\mathcal V(p,r)}{r^n} > \frac{r\xi}{\eta} > 0.
\]

In particular, $\frac{r\xi}{\eta}$ is a positive bound depending on $n$,
$g(0)$ and $T$ only. Therefore Theorem~7.2.3 allows us to pass to a subsequence
to give the following conclusion.

\begin{theorem}
\label{topping-ch8-blow-up-singularities}
\dcref{8.5.1}
\group{perelman-w-entropy}
\source{topping:page-0104}
\uses{topping-ch8-no-collapsing,topping-ch8-volume-ratio-inj-estimate,topping-ch7-compactness-ricci-flows}
Suppose that $\mathcal M$ is a closed manifold, and $g(t)$ is a Ricci flow on
a maximal time interval $[0,T)$ with $T < \infty$. Then there exist sequences
$p_i \in \mathcal M$ and $t_i \uparrow T$ with
\[
|\Rm|(p_i,t_i) = \sup_{x \in \mathcal M,\ t \in [0,t_i]} |\Rm|(x,t) \to \infty,
\]
such that, defining
\[
g_i(t) := |\Rm|(p_i,t_i)\, g\!\left(t_i + \frac{t}{|\Rm|(p_i,t_i)}\right),
\]
there exist $b = b(n) > 0$, a complete Ricci flow $(\mathcal N,\hat g(t))$
for $t \in (-\infty,b)$, and $p_\infty \in \mathcal N$ such that
\[
(\mathcal M,g_i(t),p_i) \to (\mathcal N,\hat g(t),p_\infty)
\]
as $i \to \infty$. Moreover, $|\Rm(\hat g(0))|(p_\infty)=1$, and
$|\Rm(\hat g(t))| \le 1$ for $t \le 0$.
\end{theorem}
